% archon:covers MR2223407ConstructionHilbertQuot/FlatteningStratification.lean
\chapter{Generic Flatness and the Flattening Stratification}
\label{chap:flattening}

This chapter develops the central technical device of the Quot construction. We
first prove generic flatness, in its algebraic form
(\cref{lem:generic-flatness-algebra}) and its geometric corollary
(\cref{thm:generic-flatness}). We then construct the flattening stratification of
a base scheme attached to a coherent sheaf on relative projective space
(\cref{thm:flattening-stratification}): a finite family of locally closed
subschemes, one for each fibrewise Hilbert polynomial, characterised by a
universal flattening property and with a prescribed closure order. The Quot
functor is later realised as such a stratum inside a Grassmannian.

\bigskip
\section*{Lemma on Generic Flatness}
\medskip

\begin{lemma}\leanok
[Algebraic generic flatness]
  \label{lem:generic-flatness-algebra}
  \lean{MR2223407ConstructionHilbertQuot.genericFlatness_algebra}
  % SOURCE: Nitsure, §4, lines 1711–1716 (read from references/MR2223407-construction-hilbert-quot.tex)
  % SOURCE QUOTE: "Let $A$ be a noetherian domain,
  % and $B$ a finite type $A$ algebra. Let $M$ be a finite $B$-module.
  % Then there exists an $f\in A$, $f\ne 0$, such that
  % the localisation $M_f$ is a free module over $A_f$."
  \textit{Source: Nitsure, Lemma (Generic Flatness), §4.}
  Let \(A\) be a noetherian domain and \(B\) a finite-type \(A\)-algebra. Let
  \(M\) be a finite \(B\)-module. Then there exists a nonzero \(f\in A\) such that
  the localisation \(M_f\) is a free module over \(A_f\).
\end{lemma}

% SOURCE QUOTE PROOF: "Over the quotient field $K$ of $A$, the $K$-algebra
% $B_K = K\otimes _AB$ is of finite type, and $M_K = K\otimes_AM$ is a
% finite module over $B_K$. Let $n$ be the dimension of the support
% of $M_K$ over $\spec B_K$. We argue by induction on $n$, starting
% with $n=-1$ which is the case when $M_K =0$. In this case, as
% $K\otimes _AM = S^{-1}M$ where $S = A - \{ 0\}$, each $v\in M$
% is annihilated by some non-zero element of $A$. Taking a finite
% generating set, and a common multiple of corresponding annihilating
% elements, we see there exists an $f\ne 0$ in $A$ with $fM=0$.
% Hence $M_f =0$, proving the lemma when $n=-1$.
% Now let $n\ge 0$, and let the lemma be proved for smaller values.
% As $B$ is noetherian and $M$ is assumed to be a finite $B$-module,
% there exists a finite filtration
% $$0=M_0 \subset \ldots  \subset M_r =M$$
% where each $M_i$ is a $B$-submodule of $M$ such that for each $i\ge 1$
% the quotient module $M_i/M_{i-1}$ is isomorphic to $B/\pp_i$ for
% some prime ideal $\pp_i$ in $B$.
% Note that if $0\to M'\to M\to M'' \to 0$ is a short exact
% sequence of $B$-modules, and if $f'$ and $f''$ are non-zero elements of
% $A$ such that $M'_{f'}$ and $M''_{f''}$ are free over
% respectively $A_{f'}$ and $A_{f''}$, then $M_f$ is a free module over $A_f$
% where $f=f'f''$. We will use this fact repeatedly.
% Therefore it is enough to prove the result when
% $M$ is of the form $B/\pp$ for a prime ideal $\pp$ in $B$.
% This reduces us to the case where $B$ is a domain and $M=B$.
% As by assumption $K\otimes_AB$ has dimension $n\ge 0$ (that is,
% $K\otimes_AB$ is non-zero), the map $A\to B$ must be injective.
% By Noether normalisation lemma, there exist elements
% $b_1,\ldots,b_n \in B$,
% such that $K\otimes_AB$ is finite over its subalgebra
% $K[b_1,\ldots, b_n]$ and
% the elements $b_1,\ldots, b_n$ are algebraically independent
% over $K$. (For simplicity of notation, we write $1\otimes b$ simply as $b$.)
% If $g\ne 0$ in $A$ is chosen to be a `common denominator' for
% coefficients of equations of integral dependence
% satisfied by a finite set of
% algebra generators for $K\otimes_AB$ over $K[b_1,\ldots, b_n]$,
% we see that $B_g$ is finite over $A_g[b_1,\ldots, b_n]$.
% Let $m$ be the generic rank of the finite module $B_g$ over the
% domain $A_g[b_1,\ldots, b_n]$. Then we have a short exact
% sequence of $A_g[b_1,\ldots, b_n]$-modules of the form
% $$0\to A_g[b_1,\ldots, b_n]^{\oplus m} \to B_g \to T \to 0$$
% where $T$ is a finite torsion module over $A_g[b_1,\ldots, b_n]$.
% Therefore, the dimension of the support of $K\otimes_{A_g}T$
% as a $K\otimes_{A_g}(B_g)$-module is strictly less than $n$. Hence by
% induction on $n$ (applied to the data $A_g$, $B_g$, $T$),
% there exists some $h\ne 0$ in $A$ with $T_h$ free over $A_{gh}$.
% Taking $f=gh$, the lemma follows from the
% above short exact sequence."
\begin{proof}
  Let \(K\) be the quotient field of \(A\) and put \(B_K = K\ot_A B\),
  \(M_K = K\ot_A M\); then \(B_K\) is a finite-type \(K\)-algebra and \(M_K\) is a
  finite \(B_K\)-module. Let \(n\) be the dimension of the support of \(M_K\) over
  \(\spec B_K\). We argue by induction on \(n\).

  \emph{Base case \(n=-1\)}, i.e.\ \(M_K=0\). Since \(M_K = S^{-1}M\) with
  \(S = A\setminus\{0\}\), every element of \(M\) is annihilated by a nonzero
  element of \(A\). Taking a finite generating set and a common multiple of the
  corresponding annihilators yields a nonzero \(f\in A\) with \(fM=0\), so
  \(M_f=0\), which is free over \(A_f\).

  \emph{Inductive step \(n\ge 0\).} Two reductions. First, the property
  ``\(M_f\) free over \(A_f\) for some nonzero \(f\)'' is stable under short exact
  sequences: if \(0\to M'\to M\to M''\to 0\) and \(M'_{f'}\), \(M''_{f''}\) are
  free, then \(M_f\) is free for \(f=f'f''\). As \(B\) is noetherian and \(M\)
  finite, \(M\) admits a finite filtration with successive quotients of the form
  \(B/\pp_i\); so we may assume \(M = B/\pp\), and replacing \(B\) by \(B/\pp\)
  reduces to \(B\) a domain and \(M=B\). Since \(B_K\) has dimension \(n\ge 0\) it
  is nonzero, so \(A\hookrightarrow B\). By Noether normalisation there are
  algebraically independent \(b_1,\dots,b_n\in B\) with \(B_K\) finite over
  \(K[b_1,\dots,b_n]\); clearing denominators of the integral-dependence
  equations gives a nonzero \(g\in A\) with \(B_g\) finite over
  \(A_g[b_1,\dots,b_n]\). Let \(m\) be the generic rank of \(B_g\) over the domain
  \(A_g[b_1,\dots,b_n]\), giving a short exact sequence
  \[
    0\to A_g[b_1,\dots,b_n]^{\oplus m}\to B_g \to T \to 0
  \]
  with \(T\) a finite torsion module. Then the support of \(K\ot_{A_g}T\) has
  dimension strictly less than \(n\), so by the inductive hypothesis applied to
  \((A_g, B_g, T)\) there is a nonzero \(h\in A\) with \(T_h\) free over
  \(A_{gh}\). Taking \(f=gh\), the displayed sequence shows \(M_f=B_f\) is free
  over \(A_f\).
\end{proof}

The geometric form follows by restricting to a nonempty affine open of the base.

\begin{theorem}\leanok
[Generic flatness]
  \label{thm:generic-flatness}
  \lean{MR2223407ConstructionHilbertQuot.genericFlatness}
  \uses{lem:generic-flatness-algebra}
  % SOURCE: Nitsure, §4, lines 1781–1787 (read from references/MR2223407-construction-hilbert-quot.tex)
  % SOURCE QUOTE: "Let $S$ be a noetherian and integral scheme.
  % Let $p: X\to S$ be a finite type morphism, and let $\F$ be a
  % coherent sheaf of ${\mathcal O}_X$-modules. Then there exists a non-empty
  % open subscheme $U\subset S$ such that the restriction of $\F$
  % to $X_U = p^{-1}(U)$ is flat over ${\mathcal O}_U$."
  \textit{Source: Nitsure, Theorem (Generic Flatness), §4.}
  Let \(S\) be a noetherian integral scheme, \(p:X\to S\) a finite-type morphism,
  and \(\Ff\) a coherent sheaf of \(\Oo_X\)-modules. Then there exists a nonempty
  open subscheme \(U\subset S\) such that the restriction of \(\Ff\) to
  \(X_U = p^{-1}(U)\) is flat over \(\Oo_U\).
\end{theorem}

% SOURCE QUOTE PROOF: "The above theorem has the following consequence, which follows
% by restricting attention to a non-empty affine open subscheme of $S$."
\begin{proof}
  \uses{lem:generic-flatness-algebra}
  Flatness over \(\Oo_U\) may be checked over a nonempty affine open
  \(U=\spec A\subset S\), with \(A\) a noetherian domain. Cover \(p^{-1}(U)\) by
  finitely many affine opens \(\spec B_j\) with each \(B_j\) a finite-type
  \(A\)-algebra, on which \(\Ff\) is given by a finite \(B_j\)-module \(M_j\). By
  \cref{lem:generic-flatness-algebra} there is a nonzero \(f_j\in A\) with
  \((M_j)_{f_j}\) free, hence flat, over \(A_{f_j}\). Let \(f\) be the product of
  the \(f_j\); over the nonempty open \(\spec A_f\subset U\) every \(M_j\) becomes
  flat over \(A_f\), so \(\Ff\) restricted to \(X_{\spec A_f}\) is flat over
  \(\Oo_{\spec A_f}\). Take this open as the required \(U\).
\end{proof}

\bigskip
\section*{Existence of the Flattening Stratification}
\medskip

\begin{theorem}\leanok
[Flattening stratification]
  \label{thm:flattening-stratification}
  \lean{MR2223407ConstructionHilbertQuot.flatteningStratification}
  \uses{thm:generic-flatness, def:hilbert-polynomial, thm:semicontinuity, thm:base-change-flat, lem:flatness-from-local-freeness, lem:quot-hilbert-decomposition}
  % SOURCE: Nitsure, §4, lines 1811–1841 (read from references/MR2223407-construction-hilbert-quot.tex)
  % SOURCE QUOTE: "Let $S$ be a noetherian scheme, and let
  % $\F$ be a coherent sheaf on the projective space $\PP^n_S$
  % over $S$. Then the set $I$ of Hilbert polynomials of
  % restrictions of $\F$ to fibers of $\PP^n_S\to S$
  % is a finite set. Moreover, for each $f\in I$ there exist
  % a locally closed subscheme $S_f$ of $S$,
  % such that the following conditions are satisfied.
  % \textbf{(i) Point-set: } The underlying set $|S_f|$ of $S_f$ consists of
  % all points $s\in S$ where the Hilbert polynomial of the
  % restriction of $\F$ to $\PP^n_s$ is $f$. In particular,
  % the subsets $|S_f|\subset |S|$ are disjoint,
  % and their set-theoretic union is $|S|$.
  % \textbf{(ii) Universal property: }
  % Let $S'=\coprod S_f$ be the coproduct of the $S_f$, and let
  % $i:S'\to S$ be the morphism induced by the inclusions
  % $S_f\hra S$. Then the sheaf $i^*(\F)$ on $\PP^n_{S'}$
  % is flat over $S'$. Moreover, $i:S'\to S$ has the universal property that
  % for any morphism $\varphi :T\to S$ the pullback $\varphi^*(\F)$ on
  % $\PP^n_T$ is flat over $T$ if and only if $\varphi$ factors through
  % $i:S'\to S$.  The subscheme $S_f$ is
  % uniquely determined by the polynomial $f$.
  % \textbf{(iii) Closure of strata: }
  % Let the set $I$ of Hilbert polynomials be given a total
  % ordering, defined by putting $f<g$ whenever $f(n) < g(n)$ for all
  % $n \gg  0$. Then the closure in $S$ of the subset $|S_f|$
  % is contained in the union of all $|S_g|$ where $f\le g$."
  \textit{Source: Nitsure, Theorem (Existence of Flattening Stratification), §4.}
  Let \(S\) be a noetherian scheme and let \(\Ff\) be a coherent sheaf on the
  projective space \(\PP^n_S\) over \(S\). Then the set \(I\) of Hilbert
  polynomials of the restrictions of \(\Ff\) to the fibres of \(\PP^n_S\to S\) is
  finite. Moreover, for each \(f\in I\) there is a locally closed subscheme
  \(S_f\subset S\) such that:

  \textbf{(i) Point-set.} The underlying set \(|S_f|\) consists of all \(s\in S\)
  at which the Hilbert polynomial of \(\restr{\Ff}{\PP^n_s}\) equals \(f\). In
  particular the \(|S_f|\) are pairwise disjoint and cover \(|S|\).

  \textbf{(ii) Universal property.} Let \(S'=\coprod_f S_f\) and let \(i:S'\to S\)
  be the morphism induced by the inclusions \(S_f\hookrightarrow S\). Then
  \(i^*\Ff\) on \(\PP^n_{S'}\) is flat over \(S'\), and \(i\) has the universal
  property that for any morphism \(\varphi:T\to S\) the pullback
  \(\varphi^*\Ff\) on \(\PP^n_T\) is flat over \(T\) if and only if \(\varphi\)
  factors through \(i\). The subscheme \(S_f\) is uniquely determined by \(f\).

  \textbf{(iii) Closure of strata.} Give \(I\) the total order in which \(f<g\)
  iff \(f(m)<g(m)\) for all \(m\gg 0\). Then the closure in \(S\) of \(|S_f|\) is
  contained in \(\bigcup_{f\le g}|S_g|\).
\end{theorem}

% SOURCE QUOTE PROOF: "It is enough to prove the theorem for open subschemes of
% $S$ which cover $S$, as
% the resulting strata will then glue together by their universal property.
% \textbf{Special case: } Let $n=0$, so that $\PP^n_S=S$.
% For any $s\in S$, the \textbf{fiber} $\F|_s$ of $\F$ over $s$
% will mean the pull-back of $\F$ to the subscheme
% $\spec \kappa(s)$, where $\kappa(s)$ is the residue field at $s$.
% (This is obtained by tensoring the stalk of $\F$ at $s$ with the
% residue field at $s$, both regarded as ${\mathcal O}_{S,s}$-modules.)
% The Hilbert polynomial of the
% restriction of $\F$ to the fiber over $s$ is the degree $0$
% polynomial $e \in \Q[\lambda]$, where $e = \dim_{\kappa(s)}\F|_s$.
% By Nakayama lemma, any basis of $\F|_s$ prolongs to a neighbourhood $U$ of $s$
% to give a set of generators for $F|_U$. Repeating this argument,
% we see that there exists a smaller neighbourhood $V$ of $s$ in which
% there is a right-exact sequence
% $${\mathcal O}_V^{\oplus m} \stackrel{\psi}{\to} {\mathcal O}_V^{\oplus e}
% \stackrel{\phi}{\to} \F \to 0$$
% Let $I_e\subset {\mathcal O}_V$ be the ideal sheaf
% formed by the entries of the $e\times m$ matrix
% $(\psi_{i,j})$ of the homomorphism
% ${\mathcal O}_V^{\oplus m} \stackrel{\psi}{\to} {\mathcal O}_V^{\oplus e}$.
% Let $V_e$ be the closed subscheme of $V$ defined by $I_e$.
% For any morphism of schemes
% $f: T\to V$, the pull-back sequence
% $${\mathcal O}_T^{\oplus m} \stackrel{f^*\psi}{\to} {\mathcal O}_T^{\oplus e}
% \stackrel{f^*\phi}{\to}f^* \F \to 0$$
% is exact, by right-exactness of tensor products.
% Hence the pull-back $f^*\F$ is a locally free ${\mathcal O}_T$-module of
% rank $e$ if and only if $f^*\psi =0$, that is,
% $f$ factors via the subscheme $V_e\hra V$ defined by the vanishing of all
% entries $\psi_{i,j}$.
% Thus we have proved assertions \textbf{(i)} and \textbf{(ii)} of the theorem.
% As the rank of the matrix $(\psi_{i,j})$ is lower semi-continuous,
% it follows that the function $e$ is upper semi-continuous, which
% proves the assertion \textbf{(iii)} of the theorem, completing its proof
% when $n=0$."
\begin{proof}
  \uses{thm:generic-flatness, def:hilbert-polynomial, thm:semicontinuity, lem:flatness-from-local-freeness, lem:quot-hilbert-decomposition}
  By the universal property the strata constructed over an open cover of \(S\)
  glue, so it suffices to prove the theorem locally on \(S\).

  \emph{Special case \(n=0\) (in full).} Here \(\PP^0_S=S\). For \(s\in S\) write
  \(\restr{\Ff}{s}\) for the fibre \(\Ff_s\ot_{\Oo_{S,s}}\kappa(s)\); its Hilbert
  polynomial is the constant \(e=\dim_{\kappa(s)}\restr{\Ff}{s}\). By Nakayama's
  lemma a basis of \(\restr{\Ff}{s}\) extends, over a neighbourhood \(V\) of
  \(s\), to a presentation
  \[
    \Oo_V^{\oplus m}\xrightarrow{\;\psi\;}\Oo_V^{\oplus e}
      \xrightarrow{\;\phi\;}\Ff\to 0 .
  \]
  Let \(I_e\subset\Oo_V\) be the ideal sheaf generated by the entries
  \(\psi_{i,j}\) of the \(e\times m\) matrix of \(\psi\), and let \(V_e\subset V\)
  be the closed subscheme it cuts out. For any morphism \(f:T\to V\), tensoring
  is right exact, so
  \[
    \Oo_T^{\oplus m}\xrightarrow{f^*\psi}\Oo_T^{\oplus e}
      \xrightarrow{f^*\phi} f^*\Ff\to 0
  \]
  is exact; hence \(f^*\Ff\) is locally free of rank \(e\) iff \(f^*\psi=0\),
  i.e.\ iff \(f\) factors through \(V_e\hookrightarrow V\). This proves
  \textbf{(i)} and \textbf{(ii)} for \(n=0\) (the stratum at level \(e\) is
  \(V_e\), with the stated universal property). Finally the rank of the matrix
  \((\psi_{i,j})\) is lower semicontinuous, so \(e\) is upper semicontinuous,
  which gives \textbf{(iii)}. This settles \(n=0\).

  \emph{General case (sketch).} The strategy is to realise the stratification of
  \(\Ff\) as a common refinement (a ``g.c.d.'') of the \(n=0\) flattening
  stratifications of finitely many direct images \(\pi_*\Ff(N+i)\), where
  \(\pi:\PP^n_S\to S\). The mechanism rests on \cref{lem:flatness-from-local-freeness}:
  \(\Ff\) is flat over the base precisely when the direct images
  \(\pi_*\Ff(m)\) are locally free for all large \(m\).

  \emph{Step 1 — reduction inputs.} As \(S\) is noetherian it has finitely many
  irreducible components; giving each its reduced structure and removing overlaps
  yields, via \cref{thm:generic-flatness} and noetherian induction, finitely many
  reduced, mutually disjoint, locally closed \(V_i\subset S\) covering \(|S|\)
  set-theoretically with \(\restr{\Ff}{\PP^n_{V_i}}\) flat over \(\Oo_{V_i}\).
  Because Hilbert polynomials are locally constant in a flat family
  (\cref{def:hilbert-polynomial}), this gives statement \textbf{(A)}: only
  finitely many polynomials \(P_s(m)=\chi(\PP^n_s,\Ff_s(m))\) occur as \(s\)
  ranges over \(S\); in particular the set \(I\) is finite.

  \emph{Step 2 — uniform cohomology and base change.} The semicontinuity and
  base-change machinery (\cref{thm:semicontinuity}) applied to the flat families
  over the \(V_i\) yields: \textbf{(B)} an integer \(N_1\) with
  \(R^r\pi_*\Ff(m)=0\) and \(H^r(\PP^n_s,\Ff_s(m))=0\) for all \(r\ge 1\),
  \(m\ge N_1\), \(s\in S\); and, combining the two base-change isomorphisms over
  each \(V_i\), \textbf{(C)} an integer \(N\ge N_1\) such that
  \((\pi_*\Ff(m))|_s\to H^0(\PP^n_s,\Ff_s(m))\) is an isomorphism for all
  \(m\ge N\) and \(s\in S\).

  \emph{Step 3 — assembling the strata.} Set \(E_i=\pi_*\Ff(N+i)\) for
  \(i=0,\dots,n\), coherent sheaves on \(S\). Applying the \(n=0\) case
  successively to \(E_0,E_1,\dots,E_n\) (each over the strata produced by the
  previous step) gives finitely many locally closed \(W_{e_0,\dots,e_n}\subset S\)
  through which a morphism \(f:T\to S\) factors iff every \(f^*E_i\) is locally
  free of rank \(e_i\). Via the bijection \(f\leftrightarrow(e_0,\dots,e_n)\),
  \(e_i=f(N+i)\), between numerical polynomials of degree \(\le n\) and tuples in
  \(\ZZ^{n+1}\) (here \(P_s\) has degree \(\le n\) and by \textbf{(B)} is
  determined by its values at \(N,\dots,N+n\)), relabel these as \(W_f\); then
  every \(s\in W_f\) has Hilbert polynomial \(f\). This is the decomposition of
  the base by fibrewise Hilbert polynomial recorded in
  \cref{lem:quot-hilbert-decomposition}.

  \emph{Step 4 — refining \(W_f\) to \(S_f\).} By \textbf{(C)} each
  \(\pi_*\Ff(N+i)\) has constant fibre rank \(f(N+i)\) on \(W_f\) but need not be
  locally free there; the locus of local freeness of rank \(f(N+i)\) is a closed
  subscheme \(W_f^{(i)}\subset W_f\) with \(|W_f^{(i)}|=|W_f|\), defined by a
  coherent ideal \(I_i\subset\Oo_{W_f}\) and carrying the corresponding universal
  property. By the noetherian condition the increasing chain
  \(I_0\subset I_0+I_1\subset\cdots\) stabilises, so \(I=\sum_{i\ge 0}I_i\) is
  coherent; let \(S_f\subset W_f\) be the closed subscheme it defines. Then
  \(|S_f|=|W_f|\) and every \(\pi_*\Ff(N+i)\) is locally free of rank \(f(N+i)\)
  on \(S_f\). Property \textbf{(i)} is immediate from this construction.

  \emph{Step 5 — universal property and closure.} For \textbf{(ii)}, choosing
  \(N'\ge N\) (by \cref{thm:semicontinuity}) so that base change
  \((\pi_*\Ff(m))|_{S_f}\to(\pi_{S_f})_*\Ff_{S_f}(m)\) is an isomorphism for all
  \(m\ge N'\), the direct images \(\pi_*\Ff(m)\) are locally free on \(S_f\), so
  \(\Ff_{S_f}\) is flat over \(S_f\) by \cref{lem:flatness-from-local-freeness}.
  Conversely if \(\varphi:T\to S\) has \(\Ff_T\) flat, the Hilbert polynomial is
  locally constant on \(T\); on the open-and-closed locus \(T_f\subset T\) where
  it equals \(f\) the direct images \({\pi_T}_*\Ff_T(m)\) are locally free of rank
  \(f(m)\), and this forces \(T_f\to S\) to factor schematically through
  \(S_f\) — giving the universal property and the uniqueness of \(S_f\). For
  \textbf{(iii)}, since only finitely many \(f\) occur (statement \textbf{(A)})
  there is \(p\ge N\) with \(f<g\iff f(p)<g(p)\) for all occurring \(f,g\); as
  \(S_f\) is the flattening stratum of \(\pi_*\Ff(p)\), the closure property
  follows from the upper-semicontinuity established in the \(n=0\) case applied to
  \(\pi_*\Ff(p)\). This completes the proof.
\end{proof}
